\documentclass[11pt]{article}

% ---------- Datos ----------
\newcommand{\titulo}{GeoGebra como eje de la enseñanza de la geometría analítica\\y la trigonometría en Matemáticas~I}
\newcommand{\asignatura}{Investigación Educativa Aplicada a las Matemáticas}
\newcommand{\autor}{Javier Gómez López}

% ---------- Paquetes ----------
\usepackage[utf8]{inputenc}
\usepackage[T1]{fontenc}
\usepackage[spanish,es-noshorthands]{babel}
\usepackage{lmodern}
\usepackage[paperwidth=33.867cm,paperheight=19.05cm,margin=1.1cm]{geometry}
\usepackage{xcolor}
\usepackage{enumitem}
\usepackage[most]{tcolorbox}
\usepackage{tikz}
\usetikzlibrary{arrows.meta}

\pagestyle{empty}
\setlength{\parindent}{0pt}
\renewcommand{\familydefault}{\sfdefault}

% ---------- Colores (los de la SA) ----------
\definecolor{urjc}{HTML}{C8102E}
\definecolor{rosa}{HTML}{FBE9EB}
\definecolor{titular}{HTML}{3C3C3B}
\definecolor{pistagris}{HTML}{7F7F7F}

\setlist[enumerate]{leftmargin=1.5em,itemsep=2pt,topsep=2pt,label={\color{urjc}\bfseries\arabic*.}}

\newtcolorbox{bloque}[2][]{enhanced, colback=white, colframe=urjc, boxrule=0.6pt,
  arc=2pt, left=7pt, right=7pt, top=5pt, bottom=5pt, fontupper=\normalsize,
  title={#2}, fonttitle=\Large\bfseries, colbacktitle=urjc, coltitle=white, #1}
\newtcolorbox{destacado}{enhanced, colback=rosa, colframe=rosa, arc=2pt,
  left=5pt, right=5pt, top=3pt, bottom=3pt}

% Entrada del índice: número, texto, objetivo asociado
\newcommand{\idx}[3]{\par\noindent\makebox[2.1em][l]{\color{urjc}\bfseries #1}%
  \parbox[t]{\dimexpr\linewidth-2.1em-3.2em\relax}{\raggedright #2}\hfill{\small\bfseries\color{urjc}#3}\par\vspace{2.5pt}}

\begin{document}
\color{titular}

% ---------- Cabecera: título + modelo de sombra de la SA ----------
\begin{minipage}[b]{0.72\linewidth}
{\fontsize{22}{27}\selectfont\bfseries\titulo\par}
\vspace{4pt}
{\large\color{pistagris}\autor\ \textbullet\ \asignatura\ \textbullet\ Avance del trabajo (TFM)}\par
\end{minipage}\hfill
\begin{minipage}[b]{0.27\linewidth}\raggedleft
% Adaptado de la figura «Modelo de sombra» de la situación de aprendizaje (06_reto_producto.tex)
\begin{tikzpicture}[x=0.18cm,y=0.15cm,>=Stealth,font=\footnotesize]
  \fill[pistagris!35] (-4,0) rectangle (0,12);
  \draw[thick] (-4,0) rectangle (0,12);
  \draw[thick] (-6.5,0) -- (29,0);
  \coordinate (T) at (0,12); \coordinate (S) at (24.44,0);
  \draw[urjc,thick,->] (-3.5,13.72) -- (T);
  \draw[urjc,thick] (T) -- node[above right,pos=0.42,inner sep=1pt]{$T+t\,\vec v$} (S);
  \draw[very thick,titular!80] (0,0) -- node[below,inner sep=2pt]{$L=H/\tan h$} (S);
  \draw[<->] (-5.3,0) -- node[left,inner sep=1pt]{$H$} (-5.3,12);
  \draw[very thick,blue!60!black] (25.2,0.3) -- (28.2,2.6);
  \node[blue!60!black,above,inner sep=1pt] at (27.2,2.7) {panel};
  \fill (T) circle (2pt);
  \draw (20.4,0) arc (180:154:4.04);
  \node at (19.4,1.0) {$h$};
\end{tikzpicture}\\[-1pt]
{\scriptsize\color{pistagris}El reto: modelo de la sombra}
\end{minipage}
\vspace{2pt}
{\color{urjc}\rule{\linewidth}{1.2pt}}
\vspace{6pt}

% ---------- Problema y justificación | Objetivos e hipótesis | Estructura ----------
\begin{minipage}[t]{0.315\linewidth}
\begin{bloque}[height=11.3cm]{1 \textperiodcentered{} Problema y justificación}
\begin{enumerate}
  \item \textbf{Relevancia práctica.} Vectores, trigonometría y rectas se enseñan por separado, y las carencias llegan a la universidad: en 1.º de Ingeniería fallan en trigonometría básica y no conectan los vectores con su gráfica \emph{(Mahadewsing et~al., 2024; Geewe et~al., 2024)}.
  \item \textbf{Conveniencia.} Metaanálisis recientes: GeoGebra tiene un efecto positivo, grande en geometría analítica \emph{(Juandi et~al., 2021; Alfonso-Castelló y Cabero-Fayos, 2026)}.
  \item \textbf{Valor teórico.} Apenas hay estudios de vectores con GeoGebra en Bachillerato \emph{(Aliu et~al., 2025)}.
  \item \textbf{Utilidad metodológica.} GeoGebra no funciona por sí solo: faltan criterios de diseño de tareas \emph{(Drijvers, 2015; Villalpando y Pantoja, 2022)}.
\end{enumerate}
\vspace{4pt}
\begin{destacado}
\textbf{\color{urjc}Pregunta de investigación.} ¿Qué criterios de diseño deben cumplir las tareas con GeoGebra para que el alumnado de 1.º de Bachillerato integre vectores, trigonometría y rectas y justifique sus procedimientos?
\end{destacado}
\end{bloque}
\end{minipage}\hfill
\begin{minipage}[t]{0.315\linewidth}
\begin{bloque}[height=11.3cm]{2 \textperiodcentered{} Objetivos e hipótesis}
\begin{destacado}
\textbf{\color{urjc}Objetivo general.} Diseñar una propuesta de actividades para Matemáticas~I, fundamentada en la investigación, que use GeoGebra para integrar vectores, trigonometría y rectas en un reto real.
\end{destacado}
\vspace{2pt}
\textbf{Objetivos específicos}
\begin{enumerate}[label={\color{urjc}\bfseries OE\arabic*}, leftmargin=2.5em]
  \item Revisar la investigación sobre GeoGebra y la geometría dinámica en Secundaria y Bachillerato.
  \item Identificar las dificultades descritas en trigonometría y vectores.
  \item Derivar criterios de diseño de tareas con GeoGebra desde el enfoque instrumental.
  \item Concretarlos en el reto \emph{¿dónde caben más paneles solares en nuestro instituto?} y en su evaluación.
\end{enumerate}
\vspace{4pt}
\begin{destacado}
\textbf{\color{urjc}Hipótesis de trabajo.} Si las tareas con GeoGebra se diseñan con esos criterios, el alumnado conectará mejor los tres contenidos y justificará mejor lo que hace.
\end{destacado}
\end{bloque}
\end{minipage}\hfill
\begin{minipage}[t]{0.345\linewidth}
\begin{bloque}[height=11.3cm]{3 \textperiodcentered{} Estructura del TFM}
{\color{pistagris}Propuesta didáctica \textbullet\ 40--50 págs.}\par\vspace{4pt}
\idx{1.}{\textbf{Introducción:} problema, justificación, objetivos e hipótesis}{}
\begin{tcolorbox}[enhanced, colback=rosa, colframe=urjc, boxrule=0pt, leftrule=2.5pt, arc=0pt,
  left=4pt, right=4pt, top=4pt, bottom=1pt]
\idx{2.}{\textbf{Marco teórico} (10--15 págs.)}{}
\idx{2.1}{Método de búsqueda: WoS, Scopus y Dialnet (25 fuentes)}{OE1}
\idx{2.2}{Conceptos clave}{OE1}
\idx{2.3}{Enfoques teóricos: instrumental, registros de representación, geometría dinámica}{OE1}
\idx{2.4}{Dificultades: trigonometría y vectores}{OE2}
\idx{2.5}{Experiencias previas: metaanálisis y contraevidencia}{OE1}
\idx{2.6}{Síntesis: criterios de diseño e hipótesis}{OE3}
\end{tcolorbox}
\vspace{3pt}
\idx{3.}{\textbf{Propuesta didáctica:} reto de los paneles, 20 sesiones, evaluación}{OE4}
\idx{4.}{\textbf{Conclusiones}, limitaciones y prospectiva: cómo contrastar la hipótesis en el aula}{}
{\color{pistagris}Referencias \textbullet\ Anexos}
\end{bloque}
\end{minipage}

\vfill
{\color{urjc}\rule{\linewidth}{0.6pt}}\par
\vspace{2pt}
{\small\color{pistagris}Fecha preferida de presentación: \textbf{30 de noviembre}\hfill Máster en Formación del Profesorado (Matemáticas) \textbullet\ URJC \textbullet\ Curso 2026-27}
\end{document}
